
\begin{corollaire}
\label{XIII.1.16}
Las hipótesis son las de~\Ref{XIII.1.14}. Si se supone además que $F$ es
un haz de conjuntos (\resp de ind-$\LL$-grupos) y que $f_* F$ (\resp
$\R^1 f_* F$) es constructible, entonces $f_* F$ (\resp $\R^1 f_* F$)
es localmente constante.
\end{corollaire}

El corolario resulta de~\Ref{XIII.1.14} en virtud de SGA~4 IX~2.11.

\begin{remarque}
\label{XIII.1.17}
Recordemos que la condición de que $f$ sea localmente $0$-acíclico se
satisface si $f$ es plano con fibras separables y
\ifthenelse{\boolean{orig}}
{$X$ e $Y$}
{$X$ y $S$}
son localmente noetherianos (SGA~4 XV~4.1), y que la condición de que $f$ sea localmente
$1$-asférico
\marginpar{369}
para $\LL$ se satisface si $f$ es liso, siendo $\LL$ el conjunto de los
números primos distintos de las características residuales de~$S$
(SGA~4 XV~2.1).
\end{remarque}
